\documentclass[12pt]{article}
\usepackage[utf8]{inputenc}
\usepackage[T1]{fontenc}
\usepackage[margin=1in]{geometry}
\usepackage{setspace}
\usepackage{xcolor}
\usepackage[colorlinks=true, urlcolor=blue!60!black, linkcolor=blue!60!black]{hyperref}
\setlength{\parindent}{2em}
\setlength{\parskip}{1.25em}
\renewcommand{\baselinestretch}{1.0}
\title{\textit{Reed v. Town of Gilbert}, 576 U.S. 155 (2015)}
\author{Case Brief}
\date{}
\begin{document}
\maketitle

\section{Case Metadata}
\begin{itemize}
  \item \textbf{Citation:} \textit{Reed v.\ Town of Gilbert}, 576 U.S.\ 155 (2015) (No.\ 13-502).
  \item \textbf{Argued / Decided:} Argued January 12, 2015; Decided June 18, 2015.
  \item \textbf{Vote:} 9--0 judgment reversing and remanding the Ninth Circuit; 6--3 majority opinion.
  \item \textbf{Majority Author:} Justice Clarence Thomas (joined by Chief Justice Roberts and Justices Scalia, Kennedy, Alito, and Sotomayor).
  \item \textbf{Concurrences:} Justice Alito (joined by Justices Kennedy and Sotomayor); Justice Breyer (concurring in the judgment); Justice Kagan (concurring in the judgment, joined by Justices Ginsburg and Breyer).
  \item \textbf{Dissents:} None.
  \item \textbf{Chapter / Topic:} Free Speech \& Press --- Content-Based vs.\ Content-Neutral Speech Regulations and Strict Scrutiny.
  \item \textbf{Primary Source:} \href{https://supreme.justia.com/cases/federal/us/576/155/}{Justia U.S.\ Supreme Court Center (576 U.S.\ 155)}; see also \href{https://www.law.cornell.edu/constitution-conan/amendment-1/laws-making-facial-content-based-distinctions-regarding-speech}{Cornell LII Constitution Annotated}.
\end{itemize}

\section{Facts and Procedural History}
Gilbert, Arizona enacted a comprehensive Land Development Code (``Sign Code'') prohibiting outdoor signs without a permit while exempting 23 categories of signs subject to varying rules. Three exempt categories were central:
\begin{itemize}
  \item \textbf{Ideological Signs:} Noncommercial signs communicating a message or ideas not falling into another category; permitted up to 20 square feet in all zoning districts with no durational limits.
  \item \textbf{Political Signs:} Temporary signs designed to influence the outcome of an election; permitted up to 16 square feet on residential property and 32 square feet on nonresidential or municipal property, displayed up to 60 days before a primary and 15 days after a general election.
  \item \textbf{Temporary Directional Signs Relating to a Qualifying Event:} Signs directing the public to a meeting of a religious, charitable, or nonprofit group; restricted to 6 square feet, at most four signs per property, and displayable no more than 12 hours before and 1 hour after the event.
\end{itemize}

Good News Community Church, a small congregation without a permanent building, and its pastor, Clyde Reed, placed 15 to 20 temporary signs around Gilbert early each Saturday announcing the time and location of Sunday services, removing them around midday Sunday. Gilbert's Sign Code compliance manager twice cited the Church for exceeding the 12-hour pre-event window and omitting the event date, and confiscated a sign. After accommodation efforts failed, Reed and the Church sued in the U.S.\ District Court for the District of Arizona under the First and Fourteenth Amendments. The District Court denied a preliminary injunction, and the Ninth Circuit affirmed (587 F.3d 966 (9th Cir.\ 2009)), holding that a ``cursory examination'' of a sign did not render the Code content based. On remand, the District Court granted summary judgment to Gilbert, and the Ninth Circuit again affirmed (707 F.3d 1057 (9th Cir.\ 2013)), concluding under \textit{Hill v.\ Colorado}, 530 U.S.\ 703 (2000), that the categories were content neutral because Gilbert lacked censorial motive and that the Code survived intermediate scrutiny. The Supreme Court granted certiorari and reversed.

\section{Issue(s) Presented}
\begin{enumerate}
  \item Whether a municipal sign ordinance that defines categories of signs by the type of message they convey---treating ideological, political, and temporary directional signs differently---is facially content based and therefore subject to strict scrutiny regardless of the government's benign motive or lack of viewpoint discrimination.
  \item Whether Gilbert's Sign Code survives strict scrutiny as narrowly tailored to serve compelling interests in aesthetics and traffic safety.
\end{enumerate}

\section{Holding}
Yes on the first issue; no on the second. A law that draws distinctions on its face based on the communicative content of speech---whether by subject matter, function, or purpose---is content based and triggers strict scrutiny regardless of the government's benign motive, content-neutral justification, or lack of viewpoint animus. Because Gilbert's differential restrictions on temporary directional signs were hopelessly underinclusive with respect to aesthetics and traffic safety, the Sign Code failed strict scrutiny and violated the First Amendment.

\section{Reasoning}
Writing for six Justices, Justice Thomas established a two-track framework for identifying content-based speech regulations, both triggering strict scrutiny: (1) laws that are \textbf{content based on their face} because they apply to particular speech based on the topic discussed, the idea or message expressed, or the speech's function or purpose; and (2) facially neutral laws that cannot be ``justified without reference to the content of the regulated speech'' or were adopted because of disagreement with the message (\textit{Ward v.\ Rock Against Racism}, 491 U.S.\ 781, 791 (1989)).

Applying Track 1, Thomas concluded that Gilbert's Sign Code was ``paradigmatic'' facial content discrimination: whether a sign could be 20 square feet indefinitely, 32 square feet during an election, or 6 square feet for 12 hours turned entirely on its message (e.g., an ideological quote from John Locke versus an election endorsement versus a book-club meeting announcement). Thomas then rejected three Ninth Circuit rationales:
\begin{enumerate}
  \item \textbf{Benign Motive Irrelevant at Step One:} Relying on \textit{Ward}'s purpose inquiry skipped the threshold facial question. Innocent motives do not eliminate censorship risks because future officials can weaponize facially content-based laws; \textit{Ward} applies only after a law is found facially content neutral.
  \item \textbf{Subject-Matter vs.\ Viewpoint Discrimination:} Even though the Code did not favor particular viewpoints within a category, the First Amendment's bar on content discrimination extends equally to laws singling out an entire topic or subject matter (\textit{Consolidated Edison Co.\ v.\ Public Serv.\ Comm'n}, 447 U.S.\ 530, 537 (1980)).
  \item \textbf{Speaker- and Event-Based Distinctions:} The Code was not speaker based (anyone posting a church-direction sign faced the same cap), and even genuine speaker-based distinctions trigger strict scrutiny when they reflect a content preference (\textit{Turner Broad.\ Sys., Inc.\ v.\ FCC}, 512 U.S.\ 622, 658 (1994)). Likewise, targeting signs conveying the time and place of an event regulates communicative content.
\end{enumerate}

Turning to strict scrutiny, Thomas assumed arguendo that aesthetics and traffic safety were compelling interests, but held the Code ``hopelessly underinclusive'': temporary directional signs were no greater an eyesore or traffic hazard than unlimited, larger ideological signs. Finally, Thomas emphasized that municipalities retain ample content-neutral tools---regulating sign size, materials, lighting, moving parts, portability, or postings on public property---and that narrowly tailored hazard or traffic-warning signs might survive strict scrutiny.

\section{Concurring and Dissenting Opinions}
\begin{itemize}
  \item \textbf{Justice Alito (concurring, joined by Justices Kennedy and Sotomayor):} Joined the majority in full while reassuring municipalities that \textit{Reed} leaves intact content-neutral rules governing size, location, lighting, fixed versus electronic displays, public versus private property, commercial versus residential zones, on-premises versus off-premises signs, total signs per mile, and time limits on signs advertising a ``one-time event,'' as well as government speech.
  \item \textbf{Justice Breyer (concurring in the judgment):} Joined Justice Kagan and argued that because regulatory governance inevitably distinguishes by content (e.g., securities disclosures, drug labels, safety warnings), content discrimination should function as a ``rule of thumb'' within a proportionality inquiry rather than an automatic strict-scrutiny trigger.
  \item \textbf{Justice Kagan (concurring in the judgment, joined by Justices Ginsburg and Breyer):} Warned that automatic strict scrutiny for benign subject-matter exemptions (such as address signs, ``Hidden Driveway'' warnings, or ``George Washington Slept Here'' markers) threatens thousands of sensible local ordinances and risks turning courts into a ``veritable Supreme Board of Sign Review.'' Because Gilbert's arbitrary limits failed ``even the laugh test'' under intermediate scrutiny, Kagan argued the Court should have invalidated the Code without announcing an across-the-board strict-scrutiny rule.
\end{itemize}

\section{Key Quotations}
\begin{enumerate}
  \item \textbf{Presumption Against Content-Based Laws (Thomas, J., 576 U.S.\ at 163):}
  \begin{quote}
  ``Content-based laws---those that target speech based on its communicative content---are presumptively unconstitutional and may be justified only if the government proves that they are narrowly tailored to serve compelling state interests.''
  \end{quote}
  \item \textbf{Irrelevance of Benign Motive (Thomas, J., 576 U.S.\ at 165):}
  \begin{quote}
  ``A law that is content based on its face is subject to strict scrutiny regardless of the government's benign motive, content-neutral justification, or lack of `animus toward the ideas contained' in the regulated speech.''
  \end{quote}
  \item \textbf{Subject-Matter Discrimination (Thomas, J., 576 U.S.\ at 169):}
  \begin{quote}
  ``Thus, a speech regulation targeted at specific subject matter is content based even if it does not discriminate among viewpoints within that subject matter.''
  \end{quote}
  \item \textbf{Pragmatic Warning (Kagan, J., concurring in the judgment, 576 U.S.\ at 184--85):}
  \begin{quote}
  ``The Town of Gilbert's defense of its sign ordinance---most notably, the law's distinctions between directional signs and others---does not pass strict scrutiny, or intermediate scrutiny, or even the laugh test. \dots\ This Court may soon find itself a veritable Supreme Board of Sign Review.''
  \end{quote}
\end{enumerate}

\section{Doctrinal Significance and Modern Legacy (Through 2026)}
\textit{Reed} triggered a nationwide overhaul of municipal sign codes and reshaped First Amendment doctrine far beyond signage, though the Supreme Court has since calibrated its reach:
\begin{itemize}
  \item \textbf{Robust Enforcement Against Content Exceptions:} In \href{https://supreme.justia.com/cases/federal/us/591/19-631/}{\textit{Barr v.\ American Ass'n of Political Consultants, Inc.}, 591 U.S.\ 610 (2020)}, the Court applied \textit{Reed} to invalidate a federal robocall-ban exception favoring government-debt collection speech over political speech (curing the defect by severing the exception).
  \item \textbf{Retreat from the ``Need-to-Read'' Gloss --- \textit{City of Austin} (2022):} Lower courts initially read \textit{Reed} to mean any law requiring an officer to read a sign is automatically content based. In \href{https://supreme.justia.com/cases/federal/us/596/20-1029/}{\textit{City of Austin v.\ Reagan National Advertising of Austin, LLC}, 596 U.S.\ 61 (2022)} (6--3, opinion by Justice Sotomayor), the Court rejected that ``extreme'' reading, adopting Justice Alito's \textit{Reed} concurrence to hold that an on-premises/off-premises billboard distinction is facially content neutral and subject only to intermediate scrutiny because it turns on \textit{location} rather than topic or viewpoint.
  \item \textbf{2024--2026 Refinement and Continued Bite:} The Court has carved out historically grounded exceptions while preserving \textit{Reed}'s core bite. In \href{https://supreme.justia.com/cases/federal/us/602/22-704/}{\textit{Vidal v.\ Elster}, 602 U.S.\ 286 (2024)}, the Court held that the Lanham Act's content-based (but viewpoint-neutral) ``names clause'' does not trigger \textit{Reed} strict scrutiny given trademark law's history; and in \href{https://supreme.justia.com/cases/federal/us/604/24-656/}{\textit{TikTok Inc.\ v.\ Garland}, 604 U.S.\ \_\_\_ (2025)}, the Court applied intermediate scrutiny to a foreign-adversary divestiture law justified by data-security concerns rather than content suppression. Yet in \href{https://supreme.justia.com/cases/federal/us/607/24-539/}{\textit{Chiles v.\ Salazar}, 607 U.S.\ \_\_\_ (2026)} (8--1; \href{https://www.law.cornell.edu/supremecourt/text/24-539}{Cornell LII}), Justice Gorsuch invoked \textit{Reed} to confirm that laws regulating licensed counselors' talk therapy based on communicative content or viewpoint remain presumptively unconstitutional and cannot evade searching First Amendment review by being relabeled ``professional conduct.''
\end{itemize}

\section{Pedagogical Critiques and Exam Hypotheticals}
\paragraph{Pedagogical Critiques.}
\begin{enumerate}
  \item \textbf{Overbreadth of Strict Formalism:} As Justices Kagan and Breyer emphasized, treating every facial subject-matter classification as presumptively unconstitutional either imperils benign administrative and safety rules or pressures courts to dilute strict scrutiny.
  \item \textbf{Internal Tension Between Thomas and Alito:} Justice Kagan's footnote 1 highlighted an unresolved contradiction: Alito's concurrence (needed for Thomas's 6--3 majority) labeled time limits on ``one-time event'' signs content neutral, yet Thomas's majority struck down Gilbert's rule precisely because it singled out signs stating the time and location of an event. \textit{City of Austin} ultimately resolved much of this tension in favor of Alito's more pragmatic line.
\end{enumerate}

\paragraph{Exam Hypothetical 1 (Baseline --- Fails Under \textit{Reed}).}
To reduce roadside clutter, the City of Oakridge bans all temporary lawn signs except (a) ``real estate'' signs (up to 12 sq.\ ft., permitted until sale), (b) ``civic festival'' signs (up to 8 sq.\ ft., permitted 14 days before an event), and (c) ``neighborhood garage sale'' signs (up to 4 sq.\ ft., permitted 24 hours before the sale). A local resident cited for posting a 10 sq.\ ft.\ sign advertising a weekend charity bake sale challenges the ordinance.
\begin{itemize}
  \item \textbf{Analysis:} Under \textit{Reed}, the ordinance is facially content based because an inspector must classify the sign's \textit{subject matter} (real estate vs.\ civic festival vs.\ garage sale vs.\ charity bake sale) to determine permissible size and duration. Oakridge's benign anti-clutter motive is irrelevant at Step One. Under strict scrutiny, the code fails as fatally underinclusive because a 12 sq.\ ft.\ real estate sign creates as much visual clutter as a 10 sq.\ ft.\ bake-sale sign.
\end{itemize}

\paragraph{Exam Hypothetical 2 (``Opposite Side'' Contrast --- Survives Under \textit{Reed} / \textit{City of Austin}).}
Now revise Oakridge's ordinance on the \textbf{opposite side} of the doctrinal line: Oakridge repeals all topic-based categories and instead allows any non-illuminated temporary sign up to 6 sq.\ ft.\ on private residential property for up to 30 days per year, completely bans digital billboards displaying \textit{off-premises} messages (while permitting digital \textit{on-premises} signs identifying activities occurring on the same lot), and bars all private signs on city utility poles.
\begin{itemize}
  \item \textbf{Analysis:} All three rules are facially content neutral. The size, lighting, and 30-day limits apply identically regardless of message (\textit{Reed}, Part IV; Alito concurrence). The on-/off-premises digital billboard distinction is location-based rather than topic-based under \textit{City of Austin v.\ Reagan National Advertising}, 596 U.S.\ 61 (2022), even though officials must read a sign to tell whether it refers to an on-site activity. And the evenhanded ban on utility-pole postings satisfies \textit{Taxpayers for Vincent}. Subject only to intermediate scrutiny---and assuming no illicit censorial purpose under \textit{Ward}---the revised ordinance is narrowly tailored to significant aesthetic and safety interests while leaving ample alternative channels of communication.
\end{itemize}

\end{document}
